\section{From a family of words to a shared label}\label{sec:words}

\subsection{The transversal reduction}
Let \(\W\subseteq\Sigma_1\times\cdots\times\Sigma_k\), where each
alphabet is finite. Arrange an ordered \(r\)-tuple of words as the rows
of an \(r\times k\) array. Call a column \emph{admissible} if its symbols
are all equal or pairwise distinct. If two rows coincide and all columns
are admissible, the repeated entry in each column forces that column to
be constant. Thus an array with admissible columns consists either of a
single repeated word or of \(r\) distinct words. We call \(\W\)
\emph{sunflower-free} if only the first possibility occurs.

A word \(x\) represents the transversal
\(\{(i,x_i):1\le i\le k\}\). Distinct transversals form a sunflower
precisely when every column is admissible. Indeed, a symbol occurring in
two rows of a column belongs to the common intersection and must occur
in every row of that column. Conversely, the constant columns give the
common intersection, and the other columns give disjoint petals.

To find transversals in an arbitrary \(k\)-uniform family \(\F\), color
each ground element independently and uniformly with one of \(k\) colors.
Each member receives all \(k\) colors with probability \(k!/k^k\).
Some coloring therefore retains at least \((k!/k^k)|\F|\) such members.
Reading each retained set in color order identifies it with a word, and
different coordinates use disjoint color classes. Consequently, a bound
\(M^k\) for sunflower-free word families over arbitrary alphabets gives
the set bound \((k^k/k!)M^k\le(eM)^k\). This is the transversal
reduction of Alon, Shpilka and Umans~\cite[Theorem~2.6]{ASU}, based
on the partition argument of Erd\H{o}s and Kleitman~\cite{EK}.

\subsection{The price of sharing a label}
Our starting tensor lets each of \(r\) groups choose a word independently.
Testing the columns will leave one word shared by all groups. To compare
these two tensors we measure the price of sharing one of \(|\W|\)
labels on a prescribed set of sites.

Let \(S\) be a finite set of at least two sites. Tensors belong to
\(\bigotimes_{s\in S}V_s\), with each \(V_s\) a finite-dimensional
vector space over a fixed algebraically closed field of characteristic
\(p>0\). One may take \(p=2\). Positive characteristic enters the degree
constraint in Part~\ref{part:quant}. A \emph{restriction}
\(R\longrightarrow T\) applies one linear map at each site. Tensor
products group the factors at each site. In a direct sum \(R\oplus T\),
the local spaces are direct sums at every site, and the two tensors are
added after inclusion in their respective blocks.

For a nonempty subset \(P\subseteq S\), put
\[
       E_P(j)=\bigotimes_{s\in P}e_j,\qquad
       G_P(L)=\sum_{j=1}^{L}E_P(j),
\]
with scalar factors \(1\) outside \(P\). In each summand of \(G_P(L)\),
every site in \(P\) holds the same basis label \(j\). We call this tensor
a \emph{diagonal}, or a partial diagonal when \(P\ne S\). Its basis
labels can equally well be the members of a word family of size \(L\).

A \emph{spectral point} \(\varphi\) assigns nonnegative real values to
tensors. It takes zero to zero and the scalar unit to one, is monotone
under restriction, multiplies on tensor products and adds on direct
sums~\cite{Strassen}. Every nonzero tensor restricts to the unit.
Its \emph{price} \(h(T)=\log_2\varphi(T)\) is therefore nonnegative.
Define \(w(P)=h(G_P(2))\), the price of one bit shared on \(P\). Then
\begin{equation}\label{eq:price}
       h(G_P(L))=w(P)\log_2 L,\qquad 0\le w(P)\le w(S)=1 .
\end{equation}
Indeed, tensoring diagonals multiplies their label counts, and a local
projection can discard any set of labels. Compare \(L^m\) with the
two neighboring powers of two, take prices, divide by \(m\), and let
\(m\) tend to infinity. Sending each \(e_j\) to \(1\) at a site removes
that site from the shared label, so \(w(P)\) is monotone in \(P\).
The full diagonal \(G_S(2)\) is the direct sum of two units. At a
single site the diagonal is a nonzero vector and its price is zero.

\subsection{The comparison}
A site that holds a basis label \(x\) can use it as a \emph{control}.
For any family of linear maps \(M_x\) with common domain and codomain,
the assignment
\[
             e_x\otimes v\longmapsto e_x\otimes M_xv
\]
defines a single linear map. Thus the site can read a coordinate of
its word, keep the word label, and use the symbol to choose the next
operation. The same prescription acts term by term on a sum. Operations
can write fresh basis labels as controls for later maps. Adding two
restrictions with common output spaces uses the direct sum of their
inputs. To subtract the second output, multiply its block by \(-1\)
at one site.

Partition \(S\) into \(r\) nonempty groups, and assign one symbol to
each group. For three distinct groups, the rule \((U,V)\to W\) says
that when \(U\) and \(V\) agree, \(W\) must have their common symbol.
A list of rules \emph{forces admissibility} if every assignment
satisfying them gives an admissible column. A test for a rule starts
with an auxiliary tensor \(R\). At selected sites of the three groups,
the group symbol controls a local map. If all three symbols agree, the
test must return a prescribed nonzero tensor \(Q\). Its \emph{net cost}
is \(h(R)-h(Q)\). The spaces, \(R\), and \(Q\) are fixed independently
of the values of the controls.

\begin{proposition}[Comparison of independent and shared labels]\label{prop:comparison}
Fix a spectral point and a partition of \(S\) into \(r\) groups with
\(\rho=w(S)-\sum_J w(J)>0\). Suppose a list of rules forces admissibility.
For each rule \(j\) and every finite alphabet, suppose there is a test
of net cost at most \(c_j\) that returns its fixed tensor whenever the
three symbols agree and returns zero whenever the rule is violated.
Then every nonempty sunflower-free word family of length \(k\) satisfies
\[
                 \rho\log_2|\W|\le k\sum_j c_j .
\]
\end{proposition}

\begin{proof}
Write \(N=|\W|\). The tensor \(\bigotimes_J G_J(N)\) has one summand
for every independent choice of a word in \(\W\) at each group \(J\).
Supply a separate test for each rule in each column, and let \(R\) and
\(Q\) denote the products of their input and return tensors. On each
summand, run the tests controlled by the symbols of the chosen words,
retaining the word labels.

An \(r\)-tuple other than a repeated word has an inadmissible column,
because \(\W\) is sunflower-free. Some rule there is violated. Its
zero output kills the entire summand, whatever the other tests return.
On a repeated word every test returns its prescribed tensor. Summing
over the surviving words gives the restriction
\[
        \left(\bigotimes_J G_J(N)\right)\otimes R
                 \longrightarrow G_S(N)\otimes Q .
\]
By monotonicity and Equation~\eqref{eq:price}, sharing the word across
all groups increases its price by \(\rho\log_2 N\). This increase
is at most \(h(R)-h(Q)\le k\sum_j c_j\).
\end{proof}

\subsection{Choosing the prices by a Euclidean placement}
Proposition~\ref{prop:comparison} asks for weights with a positive
surplus \(\rho\) and with bounded costs for the tests. We choose them
by placing the sites in Euclidean space. For a finite nonempty set
\(P\subset\R^d\) or its convex hull, put
\[
    \delta(P)=\sqrt{2\pi}\,\E\max_{x\in P}\langle G,x\rangle,
    \qquad G\sim N(0,I_d).
\]
This normalization of Gaussian mean width gives zero on a point and
Euclidean length on a segment, since \(\sqrt{2\pi}\,\E(G_1)_+=1\).
It is translation invariant and additive under Minkowski sums.
In the plane it is the semiperimeter of the convex hull, with a
degenerate segment counted twice in the perimeter.

\begin{theorem}[Euclidean prices {\cite[Theorem~\ext{G}{thm:main} and Section~\ext{G}{sec:spectra}]{oaG}}]\label{thm:flat}
For any field and any finite set of distinct sites \(S\subset\R^d\)
with \(D=\delta(S)>0\), there is a spectral point on tensors on \(S\)
over that field for which
\[
                   w(P)=\frac{\delta(P)}{D}
                   \qquad(\varnothing\ne P\subseteq S).
\]
\end{theorem}

We take the point induced by the compatible Euclidean assignment
of~\cite{oaG}. Thus a single evaluation prices every shared label in
the comparison. Section~\ref{sec:plane} uses semiperimeters to choose
the groups and the sites read by each rule. Section~\ref{sec:degree}
uses the joint convexity, isometry invariance and dilation law of the
same assignment to price a total-degree constraint.
